% 第三部分 C：讨论（由 main_zh.tex 引入）。

\section{讨论}\label{sec:discussion}

\paragraph{在一道题内部，评判模型几乎看不出一个解答会走向哪里。} 离线时，\kev{} 的 \emph{on track} 分数区分同一道题 128 词元前缀的正确与错误续写，AUROC 为 \NUM{aucw_128_kevtrack}，不比同样大小的 AR 评判模型（\NUM{aucw_128_artrack}）或策略自身的对数概率（\NUM{aucw_128_logprob}）更好，而后者不需要任何额外代价。合并所有题目时每种信号都更好，因为每种信号都能部分看出题目有多难，而这对在同一道题的前缀中做选择没有帮助。在训练中，面对策略自己的根，这种区分更弱（\NUM{auditotaucw_all_treekev} \ci{\NUM{auditotaucwlo_all_treekev}}{\NUM{auditotaucwhi_all_treekev}}）；保留 \kev{} 在四个根中选出的两个，让正确叶子的比例提高了 \NUM{auditgain_all_treekev} \ci{\NUM{auditgainlo_all_treekev}}{\NUM{auditgainhi_all_treekev}} 个百分点，随机选择为 \NUM{auditgain_all_treerand}，事后最好的两个根为 \NUM{auditbest_all_treekev}。这一收益在训练早期更大（前半段 \NUM{auditgain_firsthalf_treekev} 个百分点，后半段 \NUM{auditgain_secondhalf_treekev}，未给区间）；与此同时，策略的根越来越常导向正确答案，选择的余地也缩小了（事后最好的两个根：\NUM{auditbest_firsthalf_treekev} 和 \NUM{auditbest_secondhalf_treekev} 个百分点）。

\paragraph{根上的小幅收益没有改变策略学到的东西。} 在树形相同时，保留 \kev{} 选出的根而不是随机的根，最终准确率在五个基准上都在按题目重采样的区间内没有变化；GSM8K 上的差值为 \NUM{dvtreerand_treekev_gsm} \ci{\NUM{dvtreerandlo_treekev_gsm}}{\NUM{dvtreerandhi_treekev_gsm}} 个百分点，区间恰好触及零。评判模型选根背后的想法是：它能给 GRPO 更多结果有对有错的组，而只有这样的组才有学习信号。事实并非如此：在最后 20 步中，\texttt{tree\_kev} 下有 \NUM{allwrong_treekev} 的组没有任何正确叶子，\texttt{tree\_random} 下为 \NUM{allwrong_treerand}；全部叶子都正确的组分别为 \NUM{allright_treekev} 和 \NUM{allright_treerand}。这种选择在原理上有代价：被保留的根不再是策略的样本，梯度因而不再是策略梯度的无偏估计。我们没有看到它对准确率的影响。评判模型的选择也没有让策略重复自己：\texttt{tree\_kev} 的一个种子逐渐在第一个答案之后重复解答，但完全没有评判模型的 \texttt{rlvr} 也有一个种子如此，而且 \kev{} 只给解答的前 128 个词元打分，远在答案之前。

\paragraph{树形：种子均值由切换了的种子决定。} 按种子均值，两种树状条件最终都高于扁平采样：MATH-500 上分别高 \NUM{d_treerand_math} 和 \NUM{d_treekev_math} 个百分点，AIME24 上分别高 \NUM{d_treerand_aime24} 和 \NUM{d_treekev_aime24}；但这主要来自改成以纯文本开头写解答的 \NUM{treedraftruns} 个树状 run。它们在 AIME24 上最终为 \NUM{treedraftaime24min}\%--\NUM{treedraftaime24max}\%，其他树状 run 为 \NUM{treeotheraime24min}\%--\NUM{treeotheraime24max}\%，\texttt{rlvr} 为 \NUM{flataime24min}\%--\NUM{flataime24max}\%。同样的切换在整个项目中无论奖励为何都标志着最好的 run（第二部分 \S8.3）。这里树状 run 中有 \NUM{treedraftruns} 个切换，\texttt{rlvr} 的 \NUM{flatruns} 个种子一个也没有。树是否让切换更容易发生，这么少的 run 无法回答；何况这一切换不在我们的计划之中，是在 run 全部结束后才发现的，分类所用的度量也是在读过解答之后定义的。在一棵树里，解答怎样开头由根决定，所以选根同时也是在几种开头之间做选择；随机选根就足以让一个 run 发生切换。树形确定能带来的是词元上的节省：与生成同样长度解答的扁平采样相比，树状 run 每步少生成 \NUM{treegensavemin}--\NUM{treegensavemax} 的词元。

\paragraph{局限。} 只有一种策略规模（4B）、一个对 128 词元的根提出的评判问题（\emph{on track}）、一种树形（四个根，保留两个，每个四个叶子）、以第二部分的学习率训练 200 步，以及每种树状条件 \NUM{seeds_treekev} 个种子（\texttt{rlvr}：\NUM{seeds_rlvr} 个），其中第三个种子是在前两个跑完之后加的。每个提示的根更多时，评判模型可选的余地更大，它犯错的余地也更大。审计每十步才续写全部的根，每个 run 有 20 个审计步；它的区间对这些步重采样，而不对种子重采样。最终准确率的区间对题目重采样，也不对种子重采样。离线研究使用的是基座策略的前缀，而不是训练中策略的前缀。风格的度量是在读过解答之后定义的。我们报告词元数而不是时间：这些 run 与其他实验共用 GPU。

\paragraph{与早期版本的关系。} 本文第 1--3 版报告的剪枝加速、KV 缓存复用率、每节点延迟和定理证明结果（MiniF2F-Lean4、ProofNet）并未经过测量（见文首的说明和 \file{CORRECTION.md}）。这些在本文中都没有测量：我们统计词元而不是时间，也没有做定理证明。凡是我们能检验的，都检验了。v3 的剪枝规则——评判模型的 $p(\text{yes})<0.3$ 时剪掉前缀——用 \emph{on track} 问题时只剪掉 \NUM{thrpruned_128_kevtrack} 的 128 词元前缀，只抓住 \NUM{thrcaught_128_kevtrack} 的死前缀（\S\ref{sec:res-prefix}）。按 \kev{} 选择续写哪些根——这正是那些版本背后的想法——并没有让 GRPO 学得比随机选根更好。
